% =====================================================================
%  preamble.tex — gói, font, header/footer, định dạng tiêu đề
%  Bám định dạng file mẫu (Nhom_6-CS340...): Times, header 2 bên,
%  Chương N. / N.M. / N.M.K., hình/bảng đánh số toàn cục, IEEE refs.
% =====================================================================

% ---------- Ngôn ngữ & font (XeLaTeX) ----------
\usepackage{fontspec}
\usepackage[vietnamese]{babel}     % tự dịch: Hình, Bảng, Chương, Mục lục, Danh sách hình vẽ...
% Font chính: Times New Roman (có sẵn trên macOS/Windows, đủ glyph tiếng Việt).
% Linux: cài msttcorefonts hoặc đổi sang "Liberation Serif" / "TeX Gyre Termes".
\setmainfont{Times New Roman}
\setmonofont{Menlo}[Scale=0.85]   % đường dẫn, tên file, tên biến
% \setmainfont{TeX Gyre Termes}    % <- fallback đa nền tảng (cần gói tex-gyre)

% ---------- Trình bày trang ----------
\usepackage[a4paper,left=3cm,right=2cm,top=2.5cm,bottom=2.5cm]{geometry}
\usepackage{setspace}
\onehalfspacing
\usepackage{indentfirst}           % thụt lề dòng đầu mỗi đoạn (chuẩn VN)
\setlength{\parindent}{1cm}

% ---------- TikZ (UML diagrams: ER, sequence, use case, class — Ch7) ----------
\usepackage{tikz}
\usetikzlibrary{positioning,shapes.geometric,arrows.meta,fit,calc}

% ---------- Đồ họa, bảng, màu ----------
\usepackage{graphicx}
\graphicspath{{figures/}}
% Quy ước ảnh: MỌI \includegraphics viết KHÔNG ĐUÔI, ví dụ
%   \includegraphics[width=0.8\textwidth]{eda_class_distribution}
% XeLaTeX bắt .pdf (vector, nét khi in), pandoc bắt .png (nhúng được vào .docx).
% Thứ tự ưu tiên dưới đây quyết định điều đó cho phía XeLaTeX.
\DeclareGraphicsExtensions{.pdf,.png,.jpg}
\usepackage{booktabs}
\usepackage{multirow}
\usepackage{array}
\usepackage{longtable}
\usepackage{xcolor}
\definecolor{uitblue}{RGB}{0,84,159}
\usepackage[font=small,labelfont=bf]{caption}
\usepackage{enumitem}
\usepackage{amsmath,amssymb}
\usepackage{siunitx}               % số & đơn vị, dấu phẩy thập phân theo quy ước VN
\sisetup{output-decimal-marker={,}, group-separator={.}, group-minimum-digits=4}

% ---------- Hình/bảng đánh số THEO CHƯƠNG (Hình 3.1, Bảng 3.1) như BieuMau ----------
\counterwithin{figure}{chapter}
\counterwithin{table}{chapter}
\renewcommand{\thefigure}{\thechapter.\arabic{figure}}
\renewcommand{\thetable}{\thechapter.\arabic{table}}

% ---------- Header / footer 2 bên (fancyhdr) ----------
\usepackage{fancyhdr}
\pagestyle{fancy}
\fancyhf{}
\fancyhead[L]{\small UIT -- Đồ án tốt nghiệp}
\fancyhead[R]{\small \ThesisTitleShort{} -- \StudentA{} \& \StudentB}
\fancyfoot[C]{\thepage}
\renewcommand{\headrulewidth}{0.4pt}
\renewcommand{\footrulewidth}{0pt}
% Trang mở đầu chương cũng giữ header giống file mẫu:
\fancypagestyle{plain}{%
  \fancyhf{}%
  \fancyhead[L]{\small UIT -- Đồ án tốt nghiệp}%
  \fancyhead[R]{\small \ThesisTitleShort{} -- \StudentA{} \& \StudentB}%
  \fancyfoot[C]{\thepage}%
  \renewcommand{\headrulewidth}{0.4pt}%
}

% ---------- Định dạng tiêu đề: chương ngắt trang, tên IN HOA ----------
\usepackage{titlesec}
\titleformat{\chapter}[hang]
  {\normalfont\fontsize{14}{17}\bfseries}
  {\chaptertitlename\ \thechapter.}{0.6em}{\MakeUppercase}
\titlespacing*{\chapter}{0pt}{0pt}{20pt}
\titlelabel{\thetitle.\quad}       % dấu chấm sau số mục: "1.1.", "1.1.1."

% ---------- Thư mục tài liệu tham khảo IEEE ----------
\usepackage[backend=biber,style=ieee,sorting=nyt]{biblatex}
\DeclareSourcemap{\maps[datatype=bibtex]{\map{\step[fieldsource=langid, final]\step[fieldset=keywords, origfieldval]}}}
\addbibresource{refs.bib}

% ---------- Liên kết ----------
\usepackage[hidelinks]{hyperref}

% ---------- Tiện ích soạn thảo: đánh dấu phần CHƯA VIẾT ----------
% Dùng \wip{ghi chú} để chèn ghi chú "đang làm" nổi bật (dễ tìm khi rà soát).
\newcommand{\wip}[1]{\textcolor{red}{\textbf{[TODO: #1]}}}
% Dùng \ph{...} cho số liệu placeholder cần thay bằng kết quả thật.
\newcommand{\ph}[1]{\textcolor{orange}{\underline{#1}}}
